% Page 013

\seite{13}

\margmark{die ältesten Lehrer\\auf der Pfarrschule}%
an der Kirchspiels ein Kind dessen Eltern und bekannt\ob
blieben und das des Pfarrers angesehen Lauf. Kluge\unsicher\ob
eines Sünderthal (Hieß Puhl\unsicher) erfüllt als den Namen\ob
Hirsch. Es wurde späterer Schullehrer und Küster\ob
von Seffern. Auf diesen wurde ein aus Sefferweichs\ob
stammender Joh.\ Wilh. Zahmen als Lehrer und\ob
Küster vorangestellt, der bis 1821 amtirte. Ihm\ob
folgte Nikolaus Zahmen, eines Gebornen ein\ob
Pfarrwesen, war Lehrer und Küster von 1821--1833,\ob
und hinterließ nun Lehrer Joh.\ Georg Zahmen,\ob
auch ein Pfarrwesen, der vom Jahre 1833--1878
den 6.\ Sept. das Lehr- und Küsteramt treu und
fleißig vollbrachte. \darueber{(135.\ Geburtstag der Bürger [?})]

\margmark{Schule\\Sefferweich\\1843}%
Von den Filialen erhielt Sefferweich im Jahre\ob
1843 die eigene Schule mit beigegebenem, durch Lehrer\ob
Johann      geb.\ 21.2.1825   \darueber{Walderdorf}\ob
Lehrer Heinzen, aus Schönecken, Kreis Prüm,\ob
gebürtig. Dieser wirkte bis zum Jahre 1851 in\ob
Sefferweich. Außerdem war er auch an unseren anderen\ob
Hulfsschulen fleißig und verbaute in seinem Garten.

Auch nach dem zwischen Sefferweich und Seffern\ob
gelegenen Heinzenhof. Im Jahre 1851 hat\ob
der Nachfolger des Heinzen in der Person des\ob
Lehrer Matth.\ Dillenburg, gebürtig aus Idenheim und\ob
\margmark{Ält. erbaute\\Schulhaus in\\Sefferweich\\1843.}%
vorgebildet im philosophischen Seminar zu Trier, in\ob
sein Amt ein.\ob
Das gegenwärtige Schulhaus an der Dorfstraße gelegen
\pend
